\subsection{精确标签与记录界常数的一次精化}
\label{app:label-refinement}
将正文十进制参数视为精确有理数，按
$P^+=(P^{-1}+R_\gamma^{-1})^{-1}$ 及预测公式枚举，
S、U工况各32条历史均产生32个不同的后验协方差矩阵。
相等判定使用有理数运算，故精确标记与事件串的信息等价不依赖浮点去重阈值。

下面只评估标签依赖系数对现有记录界常数的影响。
对历史 $h=\gamma_{0:k}$，令 $M_t=(I-K_{\gamma_t}(P_t^-))A$，
$T_{j,h}=M_k\cdots M_{j+1}K_1(P_j^-)$（空乘积为 $I$）。偏差递推给出
$D_k=-\sum_{j:\gamma_j=1}T_{j,h}a_j$。于是取
\begin{equation}
 B(P)=b_{\rm ref}\max_{h:P_k^+(h)=P}\sum_{j:\gamma_j=1}\|T_{j,h}\|_2,
 \qquad \|D_k\|_2\le B(P_k^+).
 \label{eq:label-radius}
\end{equation}
此半径适用于全部允许策略。令 $h(P)=\lambda_{\min}(H(P))$，定义
\begin{equation}
 \begin{aligned}
 c(P)&=\frac{B(P)^2}{1-e^{-h(P)B(P)^2}}\quad(B(P)>0),\\
 F(P,Y)&=\operatorname{tr}P+c(P)[1-Y/s(P)].
 \end{aligned}
 \label{eq:label-score}
\end{equation}
$B(P)=0$ 时先取连续延拓 $c(P)=1/h(P)$。
割线界给出 $c(P)(1-e^{-D^\top H(P)D})\ge\|D\|_2^2$。
由条件静默恒等式，健康时 $\E_0F=\E_0\operatorname{tr}P$，因而
\begin{equation}
 J_k\le\E_\pi F-\E_0F\le C_{\rm label}\Delta_{\rm marked},
 \qquad C_{\rm label}=\max_{P,Y}F-\min_{P,Y}F.
 \label{eq:label-bound}
\end{equation}
最大、最小值取全部可达精确标签及 $Y\in\{0,1\}$。
若 $B(P)=0$，该标签下所有允许偏差均为零，亦可直接取 $c(P)=0$；
这是另一个有效取值，不是式\eqref{eq:label-score}的连续延拓。
表\ref{tab:label-constants}比较两种系数取值与原统一常数。
\begin{table}[H]
\centering\small
\tabcaption{标签依赖系数对记录界常数的影响}{Effect of label-dependent coefficients on the record-bound constant}
\label{tab:label-constants}
\begin{tabular}{lrrr}
\toprule
工况 & 原 $C_{\rm record}$ & 连续延拓 $C_{\rm label}$ & 零半径处取零\\
\midrule
S & 6.884171 & 6.710895 & 6.660168\\
U & 5.869119 & 5.303361 & 5.165567\\
\bottomrule
\end{tabular}
\end{table}
连续延拓取值使常数分别下降约2.52\%和9.64\%；零半径处取零时分别下降约3.25\%和11.99\%。
后一取值下，线性区间的上、下界斜率比仍约为758.14和316.09。
这些浮点评估说明，标签依赖系数仅有限改善当前常数，尚未消除上下界的明显间隙。
正文图\ref{fig:vector-budget}保留原常数，以对应主定理的统一构造。
精确标签表、两种边界取值及重建脚本见复现附件；常数小数值未作严格区间认证。
